\section{Categories}
    \begin{definition}
        A \textbf{category} $\catc$ consists of
        \begin{enumerate}
            \item a collection of objects $\ob{\catc}$
            \item for each $A,B\in\ob{\catc}$ a collection of morphisms $\homc{\catc}{A}{B}$ called the hom-set of $A$ and $B$\\
            We will write $f : A\to B$ or
$\begin{tikzcd}
	A & B
	\arrow["f", from=1-1, to=1-2]
\end{tikzcd}$ for $f\in\homc{\catc}{A}{B}$
            \item for each $A,B,C\in\ob{\catc}$ a map $\circ$ : $\homc{\catc}{A}{B}\times\homc{\catc}{B}{C}\to \homc{\catc}{A}{C}$ such that
        for morphisms $\begin{tikzcd}
            A & B & C & D
            \arrow["f", from=1-1, to=1-2]
            \arrow["g", from=1-2, to=1-3]
            \arrow["h", from=1-3, to=1-4]
        \end{tikzcd}$ we have $h\circ(g\circ f) = (h\circ g)\circ f$
            \item for each object $A\in\ob{\catc}$ we have a morphism $1_A : A\to A$ satisfying
             for all morphisms $\begin{tikzcd}
                X & A & Y
                \arrow["f", from=1-1, to=1-2]
                \arrow["g", from=1-2, to=1-3]
            \end{tikzcd}$ the identity
            \begin{equation}1_A \circ f = f \hspace{20pt} g \circ 1_A = g\end{equation}
            which we will call the identity morphism. 
        \end{enumerate}
    \end{definition}
\begin{remark}
    From $(1)$ if follows that the identity morphism $1_A$ is unique, since for $1_A,\hat{1}_A$ we have
    \[1_A = 1_A\circ \hat{1}_A = \hat{1}_A\]
\end{remark}
\begin{definition}
    A category $\catc$ is called \textbf{locally small} if all of the hom-sets in $\catc$ are sets. 
    A category $\catc$ is \textbf{small} if it is locally small and $\ob{\catc}$ is a set.
\end{definition}
\begin{definition}
    For a category $\catc$ we define the \textbf{opposite category} $\opp{\catc}$ by
    \[\ob{\opp{\catc}} = \ob{\catc}\]
    \[\homc{\opp{\catc}}{A}{B} = \homc{\catc}{B}{A}\]
    For $A,B,C\in\ob{\opp{\catc}}$ and morphisms $\begin{tikzcd}
	A & B & C
	\arrow["f", from=1-1, to=1-2]
	\arrow["g", from=1-2, to=1-3]
\end{tikzcd}$ we'll let $g\circop f = f\circ g$ and define the identity morphisms of $\opp{\catc}$ as the identity morphisms of $\catc$.
\end{definition}
The opposite category is a category since 
\[
    h\circop (g\circop f) = (g\circop f) \circ h = (f \circ g) \circ h = f\circ (g\circ h) = (g\circ h)\circop f = (h\circop g)\circop f
\]
and 
\[
    g\circop 1_A = 1_A \circ g = g \hspace{20pt} 1_A\circop f = f\circ 1_A = f
\]
\begin{definition}
    Let $\catc$ and $\catd$ be categories. A \textbf{covariant functor} $F: \catc \to \catd$ consists of 
    \begin{enumerate}
        \item a map $\ob{\catc}\to\ob{\catd}$ assigning $A \mapsto F(A)$
        \item for each $A,B\in\ob{\catc}$ a map $\homc{\catc}{A}{B}\to\homc{\catd}{F(A)}{F(B)}$ which assigns $f \mapsto F(f)$
    \end{enumerate}
    which preserves composition $F(g\circ f) = F(g)\circ F(f)$ and the identity $F(1_A) = 1_{F(A)}$.\\ We define a \textbf{contravariant functor} $F:\catc \to \catd$ to 
    be just a covariant functor $F:\opp{\catc} \to \catd$.
\end{definition}
\begin{definition}
    Let $F: \catc \to \catd$ be a functor. We call $F$ \textbf{faithful}, if it is injective on hom-sets. 
    More precisely for all $A,B\in\ob{\catc}$ and each pair of morphisms $\begin{tikzcd}
        A & B
        \arrow["g"', shift right, from=1-1, to=1-2]
        \arrow["f", shift left, from=1-1, to=1-2]
    \end{tikzcd}$ if $F(f)=F(g)$ then $f=g$.\\ 
    We say that $F$ is \textbf{full} if it is surjective on hom-sets.
    So for all $A,B\in\ob{C}$ and $g: F(A)\to F(B)$ there exists $f:A\to B$ such that $F(f)=g$.
\end{definition}
\begin{remark}
    Note that the definition doesn't say anything about the map on objects being injective or surjective,
    nor does it say that if the images of any two morphisms are equal then they are the same, since
    the functor may map the morphisms between two different objects to the same morphism in the target category.
\end{remark}
\begin{definition}
    Let $F,G: \catc \to \catd$ be covariant functors. A $\textbf{natural transformation}$ $\eta : F \to G$ is a 
    collection of morphisms in $\catd$ indexed by the objects of $\catc$. Explicitly $\eta = \{\eta_A : F(A)\to G(A) \mid A \in\ob{\catc}\}$, such that
    for all objects $A,B\in\catc$ and morphisms $f : A\to B$ the following square diagram commutes
\[\begin{tikzcd}
	{F(A)} && {F(B)} \\
	\\
	{G(A)} && {G(B)}
	\arrow["{F(f)}", from=1-1, to=1-3]
	\arrow["{\eta_A}"{description}, from=1-1, to=3-1]
	\arrow["{\eta_B}"{description}, from=1-3, to=3-3]
	\arrow["{G(f)}", from=3-1, to=3-3]
\end{tikzcd}\]
A natural transformation $\eta$ is a \textbf{natural isomorphism} if for all $A\in\ob{\catc}$ the morphism $\eta_A$ is an isomorphism.
\end{definition}